
% CREACIÓN DEL DOCUMENTO
\documentclass[
	spanish, % Idioma: spanish, english, etc.
	letterpaper, oneside
]{article}

% INFORMACIÓN DEL DOCUMENTO
\def\documenttitle {Estudio de Caso}
\def\documentsubtitle {}
\def\documentsubject { Elementos del proceso de comunicación \\ Actividad Semana 2}

\def\documentauthor {Martín Jonathan de la Cruz}
\def\coursename {Redacción en contextos virtuales}
\def\coursecode {017}

\def\universityname {Universidad Abierta y a Distancia de México}
\def\universityfaculty {Licenciatura en Derecho}
\def\universitydepartment {Redacción en contextos virtuales}
\def\universitydepartmentimage {departamentos/UnADM}
\def\universitydepartmentimagecfg {height=1.57cm}
\def\universitylocation {Roma Norte, Ciudad de México}

% INTEGRANTES, PROFESORES Y FECHAS
\def\authortable {
	\begin{tabular}{ll}
		\textbf{Alumno:}
		& \begin{tabular}[t]{l}
			 Martín Jonathan de la Cruz \\	
		\end{tabular} \\
		Matrícula: &   ES2611202040 \\
		
		\textit{Profesor:}
		& \begin{tabular}[t]{l}
			Emma Liliana \\ Padilla Cano
		\end{tabular} \\
		& \\
		\multicolumn{2}{l}{Fecha de realización: \today} \\
		% \multicolumn{2}{l}{Fecha de entrega: \today} \\
		\multicolumn{2}{l}{\universitylocation}
	\end{tabular}
}

% IMPORTACIÓN DEL TEMPLATE
\input{template}

% CITAS EN FORMATO AUTOR-AÑO (APA)
\setcitestyle{authoryear,open={(},close={)}}

% INICIO DE PÁGINAS
\begin{document}

% PORTADA
\templatePortrait

% CONFIGURACIÓN DE PÁGINA Y ENCABEZADOS
\templatePagecfg

% RESUMEN O ABSTRACT
\begin{abstractd}
	\textbf{Objetivo:} 
	Analizar los elementos del proceso comunicativo y las características de la comunicación en entornos digitales para reconocer su impacto en la interacción académica y laboral.
\end{abstractd}

% TABLA DE CONTENIDOS - ÍNDICE
\templateIndex

% CONFIGURACIONES FINALES
\templateFinalcfg

% ======================= INICIO DEL DOCUMENTO =======================

\section{Introducción}

Un mensaje académico mal formulado no solo expresa una opinión: altera la relación entre los participantes y puede convertir una diferencia de ideas en un conflicto. Ese es el problema central del caso de María y Luis. La falla no está en disentir, sino en no controlar la forma, el propósito y el efecto del mensaje dentro de un sistema comunicativo regulado por normas académicas.

El proceso comunicativo permite descomponer esa situación. De acuerdo con el enfoque revisado en el Taller de lectura y redacción I, comunicar implica relacionar emisor, receptor, mensaje, canal y contexto, además de atender la intención y las condiciones de recepción \citep{romero2022taller}. En la redacción académica, por tanto, no basta con escribir lo primero que se piensa: el mensaje debe tener claridad, coherencia y adecuación al destinatario \citep{ayala2005taller}. La reescritura cumple una función estratégica porque permite corregir el impulso inicial, ajustar el tono y reducir costos comunicativos innecesarios \citep{romero2014volver}.

\section{Análisis del caso}

En el primer momento, María actúa como emisora, Luis como receptor y el mensaje es la respuesta: ``No estoy de acuerdo con lo que dices. Creo que tu idea está mal planteada y no tiene mucho sentido. Deberías investigar mejor antes de opinar''. El canal es una plataforma digital de interacción académica, posiblemente un foro o aula virtual. El contexto es educativo, porque el intercambio ocurre entre compañeros, frente a un grupo y bajo observación docente. Cuando Luis responde, la estructura se invierte: él se convierte en emisor, María en receptora y su mensaje expresa una reacción defensiva. Después, María envía un mensaje privado al docente; ahí cambia el canal, cambia el receptor y también cambia la función del mensaje, pues ya no busca debatir con Luis, sino justificar su conducta ante una autoridad académica.

La comunicación digital intensifica el problema porque elimina señales que ayudan a modular el sentido: tono de voz, gestos, pausas y expresión facial. En un espacio presencial, esas señales pueden suavizar o precisar la intención; en un foro escrito, el mensaje queda sujeto casi por completo a las palabras elegidas. Además, el canal es público ante el grupo. Eso eleva el costo del error, porque la crítica de María no solo llega a Luis, sino que afecta su imagen frente a otros estudiantes y abre la posibilidad de que el conflicto se convierta en una toma de bandos. La frase ``así hablamos en redes sociales'' revela otra falla: María traslada un código informal a un contexto académico que exige criterio, respeto y argumentación.

El error principal de María es confundir sinceridad con pertinencia. Disentir es legítimo y necesario en la formación académica; lo problemático es hacerlo sin estructura argumentativa. Expresiones como ``tu idea está mal planteada'' y ``no tiene mucho sentido'' no identifican el punto débil del argumento, no aportan evidencia y no orientan una mejora. La frase ``deberías investigar mejor antes de opinar'' desplaza la crítica del contenido hacia la persona, por lo que Luis puede recibirla como un juicio sobre su capacidad o preparación. En términos comunicativos, aparece un ruido pragmático: la intención declarada por María no coincide con el efecto producido en Luis.

Una reformulación adecuada exige conservar el desacuerdo, pero cambiar su función. El mensaje puede quedar así: ``Luis, entiendo tu postura, aunque difiero de tu planteamiento. Considero que el argumento necesita mayor precisión, especialmente en la relación entre tus ideas y la evidencia que las sostiene. Si incorporas una fuente y explicas mejor ese vínculo, tu participación puede fortalecerse''. Esta versión no renuncia a la crítica; la vuelve operativa. Señala el problema, propone un criterio de mejora y evita atacar la identidad del compañero. El resultado es un mensaje útil para el aprendizaje, no una respuesta que obliga al receptor a defenderse.

Las implicaciones educativas y laborales son directas. En el aula virtual, un mensaje mal diseñado deteriora la confianza, inhibe la participación y desplaza el objetivo de aprendizaje hacia el conflicto interpersonal. Pierde Luis, porque se siente atacado; pierde María, porque su idea queda asociada a una forma agresiva; pierde el grupo, porque la discusión deja de centrarse en el contenido. En el ámbito laboral, el costo puede ser mayor: correos, chats institucionales y plataformas de trabajo también funcionan como canales donde una frase ambigua o descortés afecta reputación, coordinación, productividad y solución de problemas. La competencia comunicativa, por tanto, es una herramienta de control profesional: permite decidir qué decir, cómo decirlo, por qué medio y con qué efecto.

\section{Conclusión}

Mi postura es que el problema no fue el desacuerdo, sino la falta de criterio comunicativo. En un entorno académico, la opinión solo tiene valor formativo cuando se convierte en argumento, y el argumento solo funciona cuando respeta al interlocutor. María tenía derecho a cuestionar la idea de Luis, pero también tenía la responsabilidad de hacerlo con precisión, respeto y finalidad constructiva. La comunicación digital exige dominar esa decisión: escribir no para descargar una reacción, sino para producir diálogo, aprendizaje y colaboración. En la educación y en el trabajo, quien no controla su mensaje pierde capacidad de influencia; quien lo estructura con claridad gana autoridad y construye mejores relaciones.

\begin{thebibliography}{99}

\bibitem[Romero Hernández(2022)]{romero2022taller}
Romero Hernández, G. (2022). \textit{Taller de lectura y redacción I}. Klik Soluciones Educativas.

\bibitem[Ayala(2005)]{ayala2005taller}
Ayala, L. (2005). \textit{Taller de lectura y redacción: libro de consulta}. Instituto Politécnico Nacional.

\bibitem[Romero Rodríguez(2014)]{romero2014volver}
Romero Rodríguez, V. J. (2014). \textit{Volver a escribir: nueve pasos para mejorar la redacción}. Sello Editorial Javeriano.
		
\end{thebibliography}

% FIN DEL DOCUMENTO
\end{document}
